\chapter{Active Disturbance Rejection Control (ADRC)}
\label{chap:adrc}

\begin{learningobjective}
\begin{itemize}[leftmargin=*]
    \item Build a formal understanding of the ADRC architecture: tracking differentiator, extended state observer (ESO), and disturbance-compensating controller.
    \item Derive the linear ADRC (LADRC) equations for second-order plants and connect the bandwidth design rules to pole-placement theory \parencite{gao2006active}.
    \item Analyse closed-loop stability and disturbance rejection through characteristic polynomials and Lyapunov arguments as in \textcite{han2009auto}.
    \item Apply the theory to a mass--spring--damper case study with multiple bandwidth selections, supported by reproducible Python tutorials.
\end{itemize}
\end{learningobjective}

\section{Concept Map for Beginners}
\begin{conceptroadmap}
\begin{enumerate}[leftmargin=*]
    \item View ADRC as ``estimate the disturbance, then cancel it''—this story makes the extended state and observer gains feel natural.
    \item Walk through the ESO pole placement slowly; it is just classical observer design with a different target polynomial.
    \item Relate bandwidth choices to everyday terms: a faster ESO hears disturbances quickly, while controller bandwidth sets how aggressively you respond.
\end{enumerate}
\end{conceptroadmap}

\section{Concept--Math--Engineering Loop}
\begin{table}[h]
    \centering
    \small
    \renewcommand{\arraystretch}{1.2}
    \begin{tabularx}{\linewidth}{|>{\raggedright\arraybackslash}p{0.28\linewidth}|>{\raggedright\arraybackslash}p{0.34\linewidth}|>{\raggedright\arraybackslash}X|}
        \hline
        \textbf{Concept} & \textbf{Mathematics} & \textbf{Engineering Practice} \\
        \hline
        ADRC architecture & Extended state definition, disturbance lumping & Sketch block diagrams labelling measured states, estimated disturbance, and compensation path. \\
        \hline
        ESO design & Pole placement via $(s + \omega_0)^3$ & Run \texttt{tutorial\_adrc} to print \texttt{np.linalg.eigvals(A\_o - L @ C)} and confirm gains. \\
        \hline
        Bandwidth tuning & Characteristic polynomials, gain scheduling & Compare \texttt{figures/adrc\_responses.png} across bandwidth sets to see predicted transient changes. \\
        \hline
    \end{tabularx}
    \caption{How Module~\ref{chap:adrc} ties the ADRC concept to algebra and reproducible simulations.}
\end{table}

\section{Architecture Overview}
\begin{beginnerguide}
    \item[\textbf{What?}] Presents the three-block ADRC structure: tracking differentiator, extended state observer, and disturbance-compensating controller.
    \item[\textbf{Why?}] Understanding each block clarifies how ADRC achieves robustness without detailed plant models.
    \item[\textbf{When?}] Before deriving equations or tuning bandwidths for a specific plant.
    \item[\textbf{How?}] Express the plant with an extended state, design the ESO, and define the control law that cancels the disturbance estimate.
    \item[\textbf{Example}] Draw the ADRC block diagram for a second-order plant and identify where each estimated state feeds back into the control input.
\end{beginnerguide}
The ADRC philosophy treats unknown dynamics $f(t)$ and disturbances as an extended state. A canonical second-order plant is written as
\begin{equation}
    \ddot{y} = f(t, y, \dot{y}) + b u, \qquad b > 0,
    \label{eq:adrc-plant}
\end{equation}
where $b$ is the nominal control gain. ADRC comprises three blocks:
\begin{enumerate}[leftmargin=*]
    \item \textbf{Tracking Differentiator (TD)} (optional in LADRC): filters the reference signal and generates smooth derivatives.
    \item \textbf{Extended State Observer (ESO)}: estimates $y$, $\dot{y}$, and the lumped disturbance $f$.
    \item \textbf{Nonlinear or linear feedback law}: cancels the disturbance estimate and imposes desired dynamics on the remaining subsystem.
\end{enumerate}
This chapter focuses on the widely used linear ADRC variant where the TD is replaced by a simple second-order reference model.

\section{Linear ADRC Derivation}
\begin{beginnerguide}
    \item[\textbf{What?}] Step-by-step derivation of the linearised ADRC equations for second-order systems.
    \item[\textbf{Why?}] Provides explicit formulas for ESO gains and feedback coefficients needed in software implementations.
    \item[\textbf{When?}] After introducing the architecture and before coding the controller.
    \item[\textbf{How?}] Form the extended state, place observer poles via $(s+\omega_0)^3$, and design the feedback law with desired closed-loop bandwidth $\omega_c$.
    \item[\textbf{Example}] Compute $\beta_1=3\omega_0$, $\beta_2=3\omega_0^2$, $\beta_3=\omega_0^3$ for $\omega_0=200$~rad/s and verify numerically that \texttt{np.linalg.eigvals(A\_o - L @ C)} returns all poles at $-200$.
\end{beginnerguide}
Introduce the extended state $x_e = [x_1\; x_2\; x_3]^T = [y\; \dot{y}\; f]$. The ESO is designed such that the estimation error dynamics follow a third-order stable polynomial.
\subsection{Extended State Observer}
Define the observer
\begin{align}
    \dot{\hat{x}}_1 &= \hat{x}_2 - \beta_1 (\hat{x}_1 - y), \label{eq:eso1}\\
    \dot{\hat{x}}_2 &= \hat{x}_3 - \beta_2 (\hat{x}_1 - y) + b u, \label{eq:eso2}\\
    \dot{\hat{x}}_3 &= -\beta_3 (\hat{x}_1 - y), \label{eq:eso3}
\end{align}
with gains
\begin{equation}
    \beta_1 = 3 \omega_0, \qquad \beta_2 = 3 \omega_0^2, \qquad \beta_3 = \omega_0^3,
    \label{eq:eso-gains}
\end{equation}
where $\omega_0$ is the observer bandwidth. The pole placement is explicit: the error dynamics matrix is
\[
    A_o - L C =
    \begin{bmatrix}
        -\beta_1 & 1 & 0 \\
        -\beta_2 & 0 & 1 \\
        -\beta_3 & 0 & 0
    \end{bmatrix},
\]
whose characteristic polynomial is
\[
    \det(s I - (A_o - L C)) = s^3 + \beta_1 s^2 + \beta_2 s + \beta_3.
\]
Setting this polynomial equal to $(s + \omega_0)^3 = s^3 + 3 \omega_0 s^2 + 3 \omega_0^2 s + \omega_0^3$ yields the gain triplet in \eqref{eq:eso-gains}. Insert a diagnostic print in \texttt{tutorial\_adrc} to evaluate \texttt{np.linalg.eigvals(A\_o - L @ C)} and confirm numerically that all poles equal $-\omega_0$. Substituting \eqref{eq:eso1}--\eqref{eq:eso3} with these gains reveals that the error dynamics satisfy
\begin{equation}
    e^{(3)} + 3 \omega_0 e^{(2)} + 3 \omega_0^2 e^{(1)} + \omega_0^3 e = 0,
\end{equation}
analogous to placing the observer poles at $-\omega_0$. In discrete implementation we use Euler integration in the tutorial to remain consistent with control-loop timing.

\subsection{Control Law and Closed-Loop Poles}
The compensated control input is
\begin{equation}
    u = \frac{v - \hat{x}_3}{b}, \qquad v = k_p (r - \hat{x}_1) - k_d \hat{x}_2,
    \label{eq:adrc-control}
\end{equation}
with feedback gains chosen according to the desired closed-loop polynomial
\begin{equation}
    s^2 + 2 \omega_c s + \omega_c^2 = 0, \qquad k_p = \omega_c^2, \quad k_d = 2 \omega_c.
\end{equation}
Combining \eqref{eq:adrc-plant} with \eqref{eq:adrc-control} yields the characteristic polynomial
\begin{equation}
    s^2 + 2 \omega_c s + \omega_c^2 = 0,
\end{equation}
indicating that perfect disturbance cancellation reduces the plant to a double integrator with bandwidth $\omega_c$. In practice $\hat{x}_3$ approximates $f$, so robustness hinges on the ESO’s ability to track $f$ faster than the closed-loop dynamics (typically $\omega_0 \geq 3 \omega_c$).

\subsection{Code Mapping}
\begin{table}[h]
    \centering
    \small
    \renewcommand{\arraystretch}{1.2}
    \begin{tabularx}{\linewidth}{|>{\raggedright\arraybackslash}p{4cm}|X|X|}
        \hline
        \textbf{Python} & \textbf{Formula} & \textbf{Explanation} \\
        \hline
        \texttt{build\_eso()} & $\beta_1 = 3\omega_0$, $\beta_2 = 3\omega_0^2$, $\beta_3 = \omega_0^3$ & Utility in \texttt{tutorials.adrc} that generates ESO gains from the bandwidth rule. \\
        \hline
        \texttt{design\_controller()} & $k_p = \omega_c^2$, $k_d = 2\omega_c$, $u = (v - \hat{x}_3)/b$ & Implements the linear ADRC control law once $\omega_c$ and $b$ are specified. \\
        \hline
    \end{tabularx}
    \caption{Mapping LADRC formulas to the tutorial code.}
\end{table}

\section{Stability and Disturbance Rejection}
\begin{beginnerguide}
    \item[\textbf{What?}] Analyses how ESO accuracy and bandwidth choices influence closed-loop stability and disturbance suppression.
    \item[\textbf{Why?}] Ensures ADRC configurations remain stable when disturbances or modelling errors occur.
    \item[\textbf{When?}] After deriving the control law to check robustness before hardware deployment.
    \item[\textbf{How?}] Examine characteristic polynomials, Lyapunov functions, and frequency-domain interpretations.
    \item[\textbf{Example}] Show that choosing $\omega_0 \ge 3\omega_c$ maintains small estimation error and plot residual disturbances using the tutorial scripts.
\end{beginnerguide}
\subsection{Error Dynamics}
Define the estimation error $\tilde{x} = x_e - \hat{x}$. From \eqref{eq:eso1}--\eqref{eq:eso3} the error dynamics are
\begin{equation}
    \dot{\tilde{x}} = (A_o - L C) \tilde{x} + [0\; 0\; \dot{f}]^T,
\end{equation}
with
\begin{equation}
    A_o = \begin{bmatrix} 0 & 1 & 0 \\ 0 & 0 & 1 \\ 0 & 0 & 0 \end{bmatrix}, \quad L = [\beta_1\; \beta_2\; \beta_3]^T, \quad C = [1\; 0\; 0].
\end{equation}
The homogeneous part has eigenvalues at $-\omega_0$ (triple pole). If $f$ is slowly varying, $\dot{f}$ acts as a small perturbation, yielding bounded estimation error. This motivates the tuning guideline $\omega_0 \gg$ dominant frequencies of $f$.

\subsection{Lyapunov Interpretation}
Following \textcite{han2009auto}, choose the Lyapunov function $V = \tilde{x}^T P \tilde{x}$ where $P$ solves the Lyapunov equation
\begin{equation}
    (A_o - L C)^T P + P (A_o - L C) = -Q,
\end{equation}
for $Q \succ 0$. With the gain selection \eqref{eq:eso-gains}, $A_o - L C$ is Hurwitz for any $\omega_0 > 0$, guaranteeing exponential convergence of $\tilde{x}$ in the absence of disturbance derivatives. The controller in \eqref{eq:adrc-control} then renders the closed-loop error dynamics
\begin{equation}
    \ddot{e} + 2 \omega_c \dot{e} + \omega_c^2 e = \tilde{x}_3,
\end{equation}
illustrating that disturbances are attenuated according to the ESO error.

\section{Tuning Workflow}
\begin{beginnerguide}
    \item[\textbf{What?}] Practical steps for selecting bandwidths, controller gains, and noise filters in ADRC.
    \item[\textbf{Why?}] A disciplined tuning procedure accelerates deployment and avoids guesswork.
    \item[\textbf{When?}] When adapting ADRC to a new plant or retuning after requirement changes.
    \item[\textbf{How?}] Pick target $\omega_c$, set $\omega_0$, adjust gain scaling, and iterate with simulation feedback.
    \item[\textbf{Example}] Follow the provided table of slow/medium/fast bandwidths and evaluate resulting step responses to choose a balanced configuration.
\end{beginnerguide}
\begin{enumerate}[leftmargin=*]
    \item Choose the tracking bandwidth $\omega_c$ based on transient requirements (settling time roughly $4/\omega_c$).
    \item Set the observer bandwidth $\omega_0$ between three and five times $\omega_c$ to ensure fast disturbance estimation without amplifying measurement noise \parencite{gao2006active}.
    \item Validate the design via simulation, examining both tracking and control effort.
    \item If large measurement noise is present, reduce $\omega_0$ and consider filtering the measured output.
    \item For plants with uncertain $b$, identify or adaptively estimate the gain using techniques from Chapter~\ref{chap:learning}.
\end{enumerate}

\section{Case Study: Mass--Spring--Damper}
\begin{beginnerguide}
    \item[\textbf{What?}] Applies ADRC formulas to a familiar second-order plant with quantifiable performance.
    \item[\textbf{Why?}] Demonstrates how tuning selections translate into measurable overshoot, settling time, and control effort.
    \item[\textbf{When?}] After learning the general workflow—use the case as a reference implementation.
    \item[\textbf{How?}] Implement ESO and controller gains, impose actuator limits, and compare bandwidth configurations.
    \item[\textbf{Example}] Run \tutorialcmd[--output-dir figures]{adrc}\\and inspect \texttt{figures/adrc\_responses.png} to compare slow/medium/fast designs.
\end{beginnerguide}
For $m=1$, $c=0.4$, $k=4$ the nominal input gain is $b=1$. We compare three bandwidth selections:
\begin{center}
\begin{tabular}{lll}
\toprule
Scenario & $\omega_c$ & $\omega_0$ \\ \midrule
Slow & $40$ & $120$ \\
Medium & $80$ & $200$ \\
Fast & $120$ & $300$ \\
\bottomrule
\end{tabular}
\end{center}
The simulations include actuator saturation at $\pm 7$~N to highlight practical limitations. Steady-state error and percent overshoot for each case are printed by the tutorial script.

\section{Tutorial: ADRC Walkthrough}
\begin{beginnerguide}
    \item[\textbf{What?}] Script-driven exploration of ADRC responses, control effort, and observer errors.
    \item[\textbf{Why?}] Visual feedback solidifies how bandwidth choices affect tracking and disturbance rejection.
    \item[\textbf{When?}] After derivations, during tuning, or when preparing documentation.
    \item[\textbf{How?}] Execute the tutorial, generate response/control/ESO plots, and adjust parameters to see their impact.
    \item[\textbf{Example}] Modify $\omega_c$ in the script, regenerate plots, and note how control saturation increases for aggressive bandwidths.
\end{beginnerguide}
\begin{enumerate}[leftmargin=*]
    \item Run \tutorialcmd[--output-dir figures]{adrc}.\\
    \item Analyse Figure~\ref{fig:adrc-responses} to see how bandwidth choice affects transient metrics.
    \item Examine Figure~\ref{fig:adrc-controls} to understand the variation in control effort and saturation utilisation.
    \item Study Figure~\ref{fig:adrc-eso} to evaluate ESO estimation error and confirm convergence consistent with the Lyapunov discussion above.
    \item Modify parameters in \texttt{tutorial\_adrc} to explore robustness against plant parameter mismatch and measurement noise.
\end{enumerate}
See Appendix~\ref{chap:tooling} for the standard command palette that recreates these ADRC experiments and keeps figures in sync with the notes.
\IfFileExists{../figures/adrc_responses.png}{
    \begin{figure}[h]
        \centering
        \includegraphics[width=0.75\linewidth]{../figures/adrc_responses.png}
        \caption{ADRC step responses for multiple bandwidth selections.}
        \label{fig:adrc-responses}
    \end{figure}
}{
    \begin{figure}[h]
        \centering
        \fbox{\parbox{0.7\linewidth}{Run the ADRC tutorial to create \texttt{figures/adrc\_responses.png}, then recompile.}}
        \caption{Placeholder for ADRC step responses.}
        \label{fig:adrc-responses}
    \end{figure}
}
\IfFileExists{../figures/adrc_controls.png}{
    \begin{figure}[h]
        \centering
        \includegraphics[width=0.75\linewidth]{../figures/adrc_controls.png}
        \caption{Control effort comparison across ADRC bandwidths.}
        \label{fig:adrc-controls}
    \end{figure}
}{
    \begin{figure}[h]
        \centering
        \fbox{\parbox{0.7\linewidth}{Generate \texttt{figures/adrc\_controls.png} to visualise ADRC control effort.}}
        \caption{Placeholder for ADRC control effort.}
        \label{fig:adrc-controls}
    \end{figure}
}
\IfFileExists{../figures/adrc_eso_errors.png}{
    \begin{figure}[h]
        \centering
        \includegraphics[width=0.75\linewidth]{../figures/adrc_eso_errors.png}
        \caption{ESO estimation error for the medium-bandwidth design.}
        \label{fig:adrc-eso}
    \end{figure}
}{
    \begin{figure}[h]
        \centering
        \fbox{\parbox{0.7\linewidth}{Generate \texttt{figures/adrc\_eso\_errors.png} to inspect ESO convergence.}}
        \caption{Placeholder for ESO estimation error.}
        \label{fig:adrc-eso}
    \end{figure}
}

\section{Discussion and Further Reading}
\begin{beginnerguide}
    \item[\textbf{What?}] Highlights extensions, research directions, and connections to broader ADRC literature.
    \item[\textbf{Why?}] Guides further study and contextualises the linear ADRC variant within the wider field.
    \item[\textbf{When?}] After mastering the basics and seeking advanced topics such as MIMO or nonlinear ADRC.
    \item[\textbf{How?}] Review recommended papers, compare methodologies, and identify scenarios where advanced observers add value.
    \item[\textbf{Example}] Read \textcite{gao2006active} and replicate their higher-order ESO design using the provided code scaffold.
\end{beginnerguide}
The LADRC design summarised here linearises the broader nonlinear ADRC methodology in \textcite{han2009auto}. Extensions include:
\begin{itemize}[leftmargin=*]
    \item Higher-order plants and multivariable systems using generalised ESOs \parencite{gao2006active}.
    \item Integration with learning-based gain identification (Chapter~\ref{chap:learning}) to adapt to changing plant gains.
    \item Nonlinear ADRC variants where the observer uses saturation functions to improve robustness to noise.
\end{itemize}

\begin{takeaway}
\begin{itemize}[leftmargin=*]
    \item ADRC delivers disturbance rejection without detailed plant models, making it attractive to beginners once the ESO concept is internalised.
    \item Formal analysis shows that bandwidth choices correspond to explicit pole locations, providing a concrete tuning recipe.
    \item Simulation and tutorial scripts bridge the gap between closed-form derivations and practical implementation, preparing learners for advanced ADRC literature.
\end{itemize}
\end{takeaway}

\section{Module Review}
\begin{itemize}[leftmargin=*]
    \item \textbf{Architecture}: separate reference smoothing, extended state observation, and disturbance-compensating feedback.
    \item \textbf{ESO gains}: $\beta_1 = 3\omega_0$, $\beta_2 = 3\omega_0^2$, $\beta_3 = \omega_0^3$ place observer poles at $-\omega_0$.
    \item \textbf{Control law}: $u = (v - \hat{x}_3)/b$ with $k_p = \omega_c^2$, $k_d = 2\omega_c$ imposes desired second-order dynamics.
    \item \textbf{Bandwidth heuristic}: choose $\omega_0 \ge 3 \omega_c$ so disturbance estimates keep pace with the controller.
    \item \textbf{Practice}: run \tutorialcmd{adrc}\\and vary $(\omega_c, \omega_0)$ to observe trade-offs in Figures~\ref{fig:adrc-responses}--\ref{fig:adrc-eso}.
\end{itemize}

\section{Keyword Review}
\begin{vocabulary}
\begin{itemize}[leftmargin=*]
    \item \textbf{ADRC}: Active Disturbance Rejection Control; framework that estimates and cancels disturbances via an extended state.
    \item \textbf{Extended State Observer (ESO)}: observer estimating plant state and lumped disturbance with gains tied to $\omega_0$.
    \item \textbf{Bandwidth $\omega_c$, $\omega_0$}: controller and observer speed parameters guiding transient response and estimation speed.
    \item \textbf{Disturbance estimate $\hat{x}_3$}: online approximation of unknown dynamics subtracted from the input.
    \item \textbf{Tracking differentiator}: optional reference filter that provides smoothed derivatives for ADRC.
\end{itemize}
\end{vocabulary}
